\documentclass[final,5p,times]{elsarticle}

% =============================================================================
% AUTHOR CHECKLIST BEFORE SUBMISSION
% 1. Add the exact public HERA H6C DOI/file identifiers for every UVH5 input.
% 2. Add the TLE provider, acquisition timestamp, source identifier, and the
%    filtering/de-duplication rule that produced the 6364-record catalog.
% 3. [DONE 2026-07-18] Integrated-tail beam-robust set confirmed
%    (22 rows / 11 pairs at B_rel>1e-3; 9 pairs at >1e-2) and the full-set
%    calibration-residual audit completed: all 176 (pair, condition) cells
%    beam-robust retained. Reflected in Secs. 3.2, 4.6, 5, 6 and abstract.
% 4. The four-flux multiplicity bootstrap was not regenerated and is therefore
%    intentionally omitted from this manuscript.
% 5. Confirm every citation key in starlink_uemr_refs_complete.bib.
% 6. Verify the Z_PS,max definition in Sec. 2.4 (median/MAD robust z) matches
%    the exact formula in the analysis scripts before submission.
% 7. Verify whether the LOCAL-tail beam-robust count (10) is rows or unique
%    pairs; the integrated-tail count (22) is rows (= 11 unique pairs).
%    Text currently labels both as rows.
% 8. Mint an archival DOI (e.g. Zenodo) for the code/summary-product release
%    and add it to the Data and Code Availability section.
% 9. Generate figure R5 (residual-audit margins) from
%    outputs/beam_robust_residual_audit_fourflux.csv; spec in the caption.
% =============================================================================

\usepackage{kotex}
\usepackage{amsmath,amssymb,bm}
\usepackage{graphicx}
\usepackage{booktabs}
\usepackage{tabularx}
\usepackage{array}
\usepackage{siunitx}
\usepackage{xurl}
\usepackage{hyperref}
\usepackage{enumitem}
\usepackage{caption}
\usepackage{subcaption}
\usepackage{placeins}
\usepackage{stfloats}
\usepackage{makecell}

\graphicspath{{./}{./figures/}{./fourflux_figures/}{./fourflux_tail_profile_figures_local/}{./fourflux_tail_profile_figures_absint/}}
\hypersetup{colorlinks=true,linkcolor=blue,citecolor=blue,urlcolor=blue}
\setlist{nosep,leftmargin=1.5em}
\captionsetup{font=small,labelfont=bf}
\journal{Astronomy and Computing}
\sisetup{detect-all}
\emergencystretch=3em
\sloppy
\raggedbottom

\newcommand{\maybeincludegraphics}[2][]{%
  \IfFileExists{#2}{\includegraphics[#1]{#2}}{%
  \fbox{\begin{minipage}{0.84\linewidth}\centering
  Figure file not found: {\ttfamily\detokenize{#2}}
  \end{minipage}}}%
}

% -----------------------------------------------------------------------------
% Frozen result macros: four-flux/full-catalog run dated 2026-07-14
% -----------------------------------------------------------------------------
\newcommand{\NTLERecords}{6364}
\newcommand{\NCells}{27}
\newcommand{\NMainRows}{1296}
\newcommand{\NBeamPairs}{648}
\newcommand{\NNullMain}{100}
\newcommand{\NNullTail}{1000}

% Main-grid counts (N_null=100)
\newcommand{\NMainLocalStrict}{30}
\newcommand{\NMainIntStrict}{74}
\newcommand{\NMainLocalPhysEthree}{11}
\newcommand{\NMainIntPhysEthree}{57}
\newcommand{\NMainLocalPhysEtwo}{7}
\newcommand{\NMainIntPhysEtwo}{35}
\newcommand{\NMainLocalPhysEone}{1}
\newcommand{\NMainIntPhysEone}{13}
\newcommand{\NMainBeamRobustLocalEthree}{4}
\newcommand{\NMainBeamRobustIntEthree}{10}
\newcommand{\NMainFrozenOnlyLocalEthree}{1}
\newcommand{\NMainFullOnlyLocalEthree}{2}
\newcommand{\NMainFrozenOnlyIntEthree}{27}
\newcommand{\NMainFullOnlyIntEthree}{10}
\newcommand{\NMainExpectedFP}{13}

% Type-II ANOVA for log10(B_rel); eta^2 rounded for presentation
\newcommand{\AnovaRtwo}{0.7571}
\newcommand{\AnovaAdjRtwo}{0.7547}
\newcommand{\AnovaFluxEta}{0.6954}
\newcommand{\AnovaLSTEta}{0.3088}
\newcommand{\AnovaMultEta}{0.2176}
\newcommand{\AnovaBaseEta}{0.0930}
\newcommand{\AnovaBeamEta}{0.0058}
\newcommand{\AnovaMorphEta}{0.0016}
\newcommand{\AnovaInterEta}{0.0001}
\newcommand{\AnovaFluxSlope}{1.042}
\newcommand{\AnovaDropped}{48}

% N_null=1000 adaptive tail refinement
\newcommand{\NLocalTailRows}{90}
\newcommand{\NLocalTailStrict}{21}
\newcommand{\NLocalTailPhysEthree}{17}
\newcommand{\NLocalTailPhysEtwo}{9}
\newcommand{\NLocalTailBeamRobustEthree}{10}
\newcommand{\NIntTailRows}{184}
\newcommand{\NIntTailStrict}{62}
\newcommand{\NIntTailPhysEthree}{61}
\newcommand{\NIntTailPhysEtwo}{40}
\newcommand{\NIntTailBeamRobustEthree}{22}

% Synthetic ZA sweep
\newcommand{\NZAAll}{210}
\newcommand{\NZAPerBeam}{105}
\newcommand{\NZAFullLocalPhys}{0}
\newcommand{\NZAFullIntPhys}{0}
\newcommand{\NZANoBeamLocalPhys}{2}
\newcommand{\NZANoBeamIntPhys}{0}
\newcommand{\NZAFullBiasEtwo}{24}
\newcommand{\NZANoBeamBiasEtwo}{103}

% Candidate-level local-branch audit and calibration residual audit
\newcommand{\NCandidateAudit}{7}
\newcommand{\NCandidateLocalRetained}{6}
\newcommand{\NCandidateLocalDropped}{1}
\newcommand{\NCandidateBeamSensitive}{0}
\newcommand{\NCandidateBeamRobustContam}{0}
\newcommand{\NCandidateDroppedPTE}{0.01199}

% -----------------------------------------------------------------------------
% Frozen result macros: integrated-tail beam-robust residual audit, 2026-07-18
% Source: outputs/beam_robust_residual_audit_fourflux.csv (+ summary)
% -----------------------------------------------------------------------------
\newcommand{\NIntAuditRowsSel}{22}
\newcommand{\NIntAuditPairs}{11}
\newcommand{\NIntAuditPairsEtwo}{9}
\newcommand{\NIntAuditPairsMid}{2}
\newcommand{\NIntAuditConds}{16}
\newcommand{\NIntAuditCells}{176}
\newcommand{\NIntAuditEvals}{352}
\newcommand{\IntAuditMaxPTE}{0.0070}
\newcommand{\IntAuditMinBrel}{1.9\times10^{-3}}
\newcommand{\IntAuditMinBrelEtwo}{0.017}
\newcommand{\NIntAuditPairsTypical}{3}
\newcommand{\NIntAuditPairsStress}{8}

\begin{document}
\begin{frontmatter}

\title{Starlink 유사 UEMR 스트레스 주입을 위한 21-cm 지연 스펙트럼 QA:\\
국소 초과, 적분 편향 및 빔-연산자 강건성의 병렬 판정}

\author[jnu]{Yujin Kim\fnref{orcid1}}
\ead{europa5218@stu.jejunu.ac.kr}
\author[jnu]{Ji-Won Kang\corref{cor1}}
\ead{jwkang@jejunu.ac.kr}
\affiliation[jnu]{organization={Jeju National University},
                  country={Republic of Korea}}
\cortext[cor1]{Corresponding author}
\fntext[orcid1]{ORCID: \href{https://orcid.org/0009-0006-9745-2960}{0009-0006-9745-2960}}

\begin{abstract}
저궤도 위성 군집의 의도치 않은 전자기 방사(unintended
electromagnetic radiation; UEMR)는 저주파 21-cm 관측의 잠재적
간섭원이지만, 원자료 수준의 존재와 재이온화 시대(EoR) 지연
윈도우의 물리적으로 유의한 오염은 동일한 명제가 아니다. 본 연구는
Starlink 유사 UEMR 순방향 주입에 대해 국소 최댓값과 윈도우 적분
통계량을 독립적인 진단 분기로 평가하고, 경험적 PTE, 상대 절대편향
$B_{\rm rel}$, 짝지은 빔 연산자 및 보정 잔차를 결합한 계산적 QA
판정 계층을 제안한다.

2022년 10월 8일 HERA H6C 파생 가시성의 시간--주파수 구조를
2026년 1월 1일 UTC 위성 기하 창으로 재매핑하고, 전체
\NTLERecords{}개 TLE 기록을 사용하였다. 9개 기선과 세 LST 층으로
구성한 \NCells{}개 셀에서 frozen/full-chromatic HERA H2C CST
PolyBeam, 세 스펙트럼 형태, 단일/다중 위성 및
$S_{\rm ref}=\{10,30,300,1000\}$ Jy를 교차하여 \NMainRows{}개
시행을 계산하였다. 거친 $N_{\rm null}=\NNullMain$ 그리드에서는
엄격한 국소 행 \NMainLocalStrict{}건과 적분 행 \NMainIntStrict{}건이
나타났으나, $B_{\rm rel}>10^{-3}$을 함께 요구하면 각각
\NMainLocalPhysEthree{}건과 \NMainIntPhysEthree{}건으로 줄었다.
Type-II 분산 분해에서 $\log_{10}B_{\rm rel}$은 플럭스
($\eta_p^2=\AnovaFluxEta$), LST 층($\AnovaLSTEta$), 위성 다중성
($\AnovaMultEta$) 및 기선 길이($\AnovaBaseEta$)가 지배했으며,
형태학과 빔의 효과크기는 작았다.

$N_{\rm null}=\NNullTail$ 적응 꼬리 정밀화에서는
$B_{\rm rel}>10^{-3}$ 기준으로 빔-강건 국소 QA 행
\NLocalTailBeamRobustEthree{}건과 빔-강건 적분 QA 행
\NIntTailBeamRobustEthree{}건(짝지은 고유 쌍 \NIntAuditPairs{}개)이
확인되었다. 이 \NIntAuditPairs{}쌍 전체에 대해 짝지은 양쪽 빔과
백색·매끄러운 보정 잔차 \NIntAuditConds{}개 조건을 교차한 전수
잔차 감사를 수행한 결과, \NIntAuditCells{}개 (쌍, 조건) 셀 전부에서
빔-강건 적분 분류가 유지되었다(최대 $PTE_{\rm global,absint}$
\IntAuditMaxPTE{}, 최소 $B_{\rm rel}=\IntAuditMinBrel$). 따라서
이들은 본 QA 계층의 최종 단계인 오염 후보 승격 기준을 만족한다.
이는 주입된 스트레스 신호에 대한 판정 계층의 종단 안정성 검증이며,
실제 우주론적 오염의 검출이 아니다. 별도의 210-case 천정각
스윕에서는 full PolyBeam 조건의 물리적 국소·적분 후보가 모두
0건이었고, 거친 국소 후보 \NCandidateAudit{}건의 독립 잔차
감사에서도 빔-강건 적분 오염 후보가 나타나지 않았다. 본 결과는
통계적 희귀성, 절대 효과크기 및 분석 연산자 의존성을 분리해
보고해야 함을 보여 주며, 실제 HERA--Starlink 동시 검출이나 발생률
추정이 아닌 재현 가능한 pre-production QA 연구로 해석되어야 한다.
\end{abstract}

\begin{keyword}
Starlink \sep unintended electromagnetic radiation \sep HERA \sep
21-cm cosmology \sep delay spectrum \sep radio-frequency interference
\sep empirical PTE \sep quality assurance \sep PolyBeam
\end{keyword}

\end{frontmatter}

% =============================================================================
\section{서론}
\label{sec:intro}

적색편이 21-cm 신호는 초기 우주의 중성수소 분포와 첫 천체 형성을
추적하는 직접적 탐침이지만, 전경, 열잡음, 보정 잔차 및 RFI보다
훨씬 희미하다
\citep{furlanetto2006,pritchard2012,morales2010,deboer2017,
abdurashidova2022}. 지연 스펙트럼 분석에서는 주파수에 대해
매끄러운 전경이 주로 지평선 지연 안쪽의 웻지에 집중되고, 지평선
버퍼 바깥의 모드를 EoR 윈도우로 사용한다
\citep{parsons2012delay,pober2013,liu2014wedge,thyagarajan2015}.
따라서 시간--주파수 평면에서 약한 구조라도 처리 연산자를 거쳐
고지연 모드로 재조직되면 과학 산출물의 QA 문제가 될 수 있다.

저궤도 위성 군집은 이 문제를 구체화한다. LOFAR와 EDA2 관측은
Starlink 위성의 저주파 의도적·비의도적 방출을 검출했고, 후속
관측에서는 광대역 및 빗살(comb-like) 구조가 보고되었다
\citep{divruno2023,grigg2023,bassa2024}. 그러나 UEMR의 존재,
특정 지연 빈의 초과, 윈도우 전체의 적분 편향, 그리고 그 판정의
빔·보정 조건에 대한 강건성은 구분되어야 한다.

이 구분은 제거 알고리즘의 성능만으로 해결되지 않는다.
\citet{kim2026nonidentifiability}은 단일 epoch에서 과학 신호와
구조화된 RFI가 중첩된 특이벡터 부분공간을 점유하면 저랭크 분해의
rank 선택만으로 잔류 RFI 억제와 신호 보존을 동시에 인증할 수
없음을 보였다. 본 연구는 이 운영적 비식별성의 문제를 처리 후 QA로
확장한다. 중심 질문은 ``RFI가 존재하는가''가 아니라, ``관측된
통계적 초과를 어느 단계까지 후속 검토가 필요한 물리적 QA 후보로
승격할 수 있는가''이다.

국소 최댓값과 윈도우 적분량은 서로 다른 신호 조직화에 반응한다.
좁은 결맞음 스파이크는 국소 통계량만 증가시킬 수 있고, 여러 지연
모드에 분산된 약한 구조는 어느 한 빈에서도 극단적이지 않으면서
적분량을 증가시킬 수 있다. 따라서 본 연구는 두 통계량을 병렬
분기로 계산하고, 통계적 희귀성, 절대 효과크기, 빔 연산자 교체 및
보정 잔차에 대한 안정성을 별도 축으로 평가한다.

본 연구는 실제 HERA--Starlink 동시 관측, 결정론적 UEMR 감산,
보정된 우주론적 밴드파워 또는 오염 발생률을 주장하지 않는다.
TLE는 개별 위성의 실제 방출 파형을 복원하는 템플릿이 아니라,
물리적으로 가능한 이동 방출원의 방향, 거리 및 지연 궤적을 생성하는
스트레스 입력이다. 본 연구의 기여는 (i) 국소·적분 반응의 병렬
판정, (ii) 문헌 규모에서 상한 스트레스까지 연결하는 4단계 플럭스
축, (iii) HERA H2C CST PolyBeam의 주파수 색수차를 포함한 짝지은
빔 감사, (iv) 빔-강건 적분 후보 전수에 대한 보정 잔차 감사,
(v) 현재 QA 연산자로 재실행한 천정각 감도 분석이다.

% =============================================================================
\section{자료, 순방향 주입 및 QA 지표}
\label{sec:methods}

\subsection{HERA 배경과 시간 재매핑}
\label{sec:background}

분석 배경은 HERA H6C 공개 가시성 산출물에서 추출한 9개 동서 방향
유사 기선의 10분 크롭이다. 원자료의 관측일은 2022년 10월 8일이며,
시간--주파수 구조, 플래그 및 \texttt{nsamples} 가중치는 그대로
유지한다. 위성 기하 계산에서는 각 크롭의 내부 표본 간격을 유지한
채 시간축만 2026년 1월 1일 UTC의 대응 창으로 재매핑하였다. 따라서
본문의 2026년 1월 1일은 TLE 평가 기준일이며 HERA 가시성의 실제
관측일이 아니다.

플래그된 표본은 가중치 0으로 두고, 비플래그 표본은
\texttt{nsamples}에 비례한 뒤 크롭별 최대값으로 정규화한다. 비유한
복소 가시성은 0으로 대체하고 해당 가중치도 0으로 둔다. 이 재매핑은
실제 HERA 관측과 Starlink 통과 사이의 동시대 위상 정합을 재현하지
않는다. 따라서 이후 경험적 PTE는 UEMR 부재 영가설의 검정값이 아니라,
고정된 배경 실현 아래에서 특정 위상 결맞음 구조가 QA 통계량을 얼마나
움직였는지를 나타내는 조건부 민감도 지표이다.

9개 기선은 14.6--207.3 m 범위를 포괄한다. 각 기선에서 주입 전
메타데이터만 사용해 정온(quiet), 전형(typical), 스트레스(stress)
LST 층을 하나씩 선택하여 총 \NCells{}개 기선--LST 셀을 구성한다.
층 선택에는 주입 후 $Z_{\rm PS,max}$, PTE 또는 $B_{\rm rel}$을
사용하지 않는다.

\subsection{TLE 기반 순방향 주입과 빔 모델}
\label{sec:injection}

전체 \NTLERecords{}개 TLE 기록을 공통 입력으로 사용하고, 각 10분
창에서 고도 기준을 통과한 위성만 선택한다. 단일 위성 조건은 가장
높은 통과 위성 1개를, 다중 위성 조건은 최대 12개를 사용한다. TLE는
결정론적 감산 모델이 아니라 시간 종속 방향·거리·지연·프린지 레이트
궤적을 생성하는 입력이다. 안테나 쌍 $(i,j)$의 근거리 지연은
\begin{equation}
 \tau_{ij}(t)=
 \frac{|\mathbf r_{\rm sat}(t)-\mathbf r_j|-|\mathbf r_{\rm sat}(t)-\mathbf r_i|}{c}
\end{equation}
로 계산한다. 위성 성분의 겉보기 진폭은
\begin{equation}
 A(t,\nu)=S_{\rm ref}M_{\rm spec}(\nu)B(t,\nu)R(t)A_{\rm smear}(t,\nu)
\end{equation}
이며, 주입 가시성은
\begin{equation}
 V_{\rm sat}(t,\nu)=A(t,\nu)\exp[-2\pi i\nu\tau_{\rm sat}(t)]
\end{equation}
이다. 다중 위성 기본 주입은 각 성분을 동일한 위상 기준에서 합한
coherent-default 상한 스트레스 구성으로 둔다. 이는 실제 서로 다른
위성 방출원의 위상 동기를 예측한다는 뜻이 아니다.

빔 응답 $B(t,\nu)$에는 \texttt{hera\_sim} v4.3의 HERA H2C CST
기반 PolyBeam을 사용한다
\citep{fagnoni2021,wilensky2024beam}. 주파수별 CST 적합 패턴을
평가하는 full-chromatic PolyBeam을 주 물리 모형으로 두고, 동일한
150 MHz 패턴을 전 대역에 고정한 frozen PolyBeam을 짝지은 연산자
대조군으로 사용한다. CST 기반 PolyBeam은 단순 Airy disk 또는
Gaussian 빔과 달리 전자기적 색수차와 비-Bessel형 null 구조를
부분적으로 반영하므로, 지연 변환 뒤 기본 위성 지연 주변의 파워를
재분배하고 해석적 빔에서 약한 고지연 응답을 만들 수 있다. 다만 이
모형은 mutual coupling, cross-polarization leakage, cable-reflection
residual 및 배열 수준의 모든 계통효과를 완전하게 포함하지 않는다
\citep{barry2016,ewallwice2017,dillon2020redcal,kim2023beamvar}.

\subsection{플럭스 계층과 스펙트럼 형태}
\label{sec:flux}

$S_{\rm ref}$는 특정 Starlink 통과의 HERA-apparent flux를 보정한
예측값이 아니라 문헌 기반 order-of-magnitude 스트레스 좌표이다.
LOFAR 관측은 저주파 Starlink UEMR이 대략 Jy에서 수십 Jy 규모에
걸쳐 나타날 수 있음을 보여 준다
\citep{divruno2023,bassa2024}. 빔 결합, slant range, 편광,
off-axis 이득 및 방출 방향성은 실제 겉보기 진폭을 크게 바꿀 수
있으므로, 본 연구는 10 및 30 Jy를 문헌 규모의 low/high anchor,
300 및 1000 Jy를 QA 연산자의 동적범위와 실패 경계를 탐색하기 위한
의도적 상한 스트레스 계층으로 사용한다. 300--1000 Jy 조건은 발생률
또는 실제 검출 플럭스에 대한 주장이 아니다.

스펙트럼 형태는 (i) 매끄러운 광대역(smooth), (ii) 문헌 기반
협대역 선형(lines), (iii) 문헌 기반 kHz-comb 뱅크로 구성한다.
형태학의 영향은 국소 PTE와 적분 편향에서 별도로 평가한다.

\subsection{지연 스펙트럼 통계량과 대응-영가설}
\label{sec:metrics}

각 10분 가시성 크롭은 짝수·홀수 적분 집합 $a,b$로 나누고,
주파수 가중치 $w(t,\nu)$와 Blackman--Harris taper $h(\nu)$를
적용한다.
\begin{equation}
 \widetilde V_{a,b}(\tau)=
 \mathcal F_{\nu\rightarrow\tau}
 \left[w_{a,b}(t,\nu)h(\nu)V_{a,b}(t,\nu)\right].
\end{equation}
별도의 spectral inpainting과 fringe-rate filtering은 적용하지 않는다.
짝수·홀수 교차곱의 실수부를 시간 평균한 QA 파워 프록시는
\begin{equation}
 \widehat P_{\rm QA}(\tau)=
 \left\langle\mathrm{Re}
 [\widetilde V_a(\tau)\widetilde V_b^*(\tau)]\right\rangle_t,
\end{equation}
주입 전후 차이는
\begin{equation}
 \Delta\widehat P_{\rm QA}(\tau)=
 \widehat P_{\rm QA}^{\rm inj}(\tau)-\widehat P_{\rm QA}^{\rm bg}(\tau)
\end{equation}
로 둔다.

EoR 윈도우 프록시는
\begin{equation}
 \tau_{\rm win}=\tau_{\rm hor}+100~{\rm ns},\qquad
 W_{\rm win}=\{\tau:|\tau|\ge\tau_{\rm win}\}
\end{equation}
로 정의한다
\citep{parsons2012delay,pober2013,liu2014wedge}. 국소 분기와 적분
분기의 시행 수준 통계량은 각각
\begin{align}
 T_{\rm max}&=\max_{\tau\in W_{\rm win}}
              \Delta\widehat P_{\rm QA}(\tau),\\
 T_{\rm absint}&=\sum_{\tau\in W_{\rm win}}
              |\Delta\widehat P_{\rm QA}(\tau)|
\end{align}
이다. 상대 절대편향은
\begin{equation}
 B_{\rm rel}=\frac{T_{\rm absint}}
 {\sum_{\tau\in W_{\rm win}}|\widehat P_{\rm QA}^{\rm bg}(\tau)|}
\end{equation}
로 정의한다. 이는 보정된 우주론적 밴드파워가 아니라 본 QA 연산자
내의 무차원 효과크기이다.

대응-영가설은 주입의 시간--주파수 지지 영역과 진폭 포락선을
유지한 채 지정된 위상 결맞음만 무작위화한다. 시행 통계량 $T$의
상측 경험적 PTE는
\begin{equation}
 PTE(T_{\rm obs})=
 \frac{1+\#\{T_{\rm null}^{(r)}\ge T_{\rm obs}\}}{N_{\rm null}+1}
\end{equation}
로 계산한다. $PTE_{\rm global,max}$는 $T_{\rm max}$,
$PTE_{\rm global,absint}$는 $T_{\rm absint}$에 대응한다. 여기서
``global''은 한 시행의 지연 윈도우 내부 최댓값 또는 적분에 대한
뜻이며, \NMainRows{}개 시행 전체의 가족 단위 보정을 뜻하지 않는다.

보조 탐색 진단으로, 대응-영가설 앙상블에 대한 윈도우 최댓값의
강건 z-점수
\begin{equation}
 Z_{\rm PS,max}=
 \frac{T_{\rm max}-\mathrm{median}\!\left(T_{\rm null}^{(r)}\right)}
      {1.4826\,\mathrm{MAD}\!\left(T_{\rm null}^{(r)}\right)}
\end{equation}
를 정의한다. $Z_{\rm PS,max}$는 꼬리 선택
(Section~\ref{sec:design:tail})과 탐색적 시각화
(Fig.~\ref{fig:z-pte})에만 사용하며, 판정 게이트로는 사용하지
않는다.

\subsection{병렬 분기와 후보 용어}
\label{sec:hierarchy}

Table~\ref{tab:hierarchy}와 같이 국소 분기와 적분 분기를 독립적으로
평가한다. $B_{\rm rel}>10^{-3}$은 탐색적 효과크기 게이트,
$B_{\rm rel}>10^{-2}$는 운영적 보고 게이트로 사용한다. 어느 값도
실제 우주론적 오차 예산으로 보정된 보편 임계값은 아니다.

본 논문에서 ``빔-강건''은 동일한 (기선, LST 층, 형태학, 플럭스,
다중성) 조건에서 frozen과 full-chromatic PolyBeam이 각각
독립적으로 해당 분기의 PTE 게이트와 $B_{\rm rel}$ 게이트를
통과함을 뜻하며, 주 그리드·꼬리 정밀화·잔차 감사 전체에서 동일한
정의를 사용한다. 표의 마지막 단계 ``오염 후보 승격''은 주입
시행에 대한 QA 계층 내 등급이다. 실제 관측 파이프라인에서라면
승격된 후보는 심층 감사를 요구하는 등급에 해당하지만, 본 연구의
주입은 합성 스트레스이므로 승격은 오염의 검출이 아니라 판정 계층의
종단(end-to-end) 안정성 검증으로 해석된다.

\begin{table}[!t]
\centering
\small
\caption{병렬 QA 판정 규칙. ``빔-강건''은 동일 물리 조건의 frozen 및
full-chromatic PolyBeam 양쪽에서 해당 분기와 동일한
$B_{\rm rel}$ 게이트가 유지됨을 뜻한다.}
\label{tab:hierarchy}
\begin{tabularx}{\linewidth}{lX}
\toprule
단계 & 규칙 및 해석 \\
\midrule
국소 분기 L & $PTE_{\rm global,max}<0.01$ \\
적분 분기 I & $PTE_{\rm global,absint}<0.01$ \\
효과크기 & $B_{\rm rel}>10^{-3}$ 또는 $10^{-2}$ \\
빔-강건 국소 QA & L과 효과크기가 양쪽 PolyBeam에서 유지 \\
빔-강건 적분 QA & I와 효과크기가 양쪽 PolyBeam에서 유지 \\
오염 후보 승격 & 빔-강건 적분 QA가 후보 수준 잔차 감사에서도 유지 \\
\bottomrule
\end{tabularx}
\end{table}

% =============================================================================
\section{실험 설계}
\label{sec:design}

\subsection{전체 카탈로그 4-플럭스 커버리지 그리드}
\label{sec:design:grid}

전체 \NTLERecords{}개 TLE를 모든 시행에 공통 적용한다. 9개 기선
$\times$ 세 LST 층으로 구성한 \NCells{}개 셀에 두 빔, 세 형태학,
네 플럭스 및 두 위성 다중성 수준을 교차한다.
\begin{equation}
 27\times2\times3\times4\times2=\NMainRows{}.
\end{equation}
각 행은 $N_{\rm null}=\NNullMain$로 계산하며, 동일한 셀·형태학·플럭스·
다중성의 frozen/full PolyBeam을 짝지어 \NBeamPairs{}쌍을 구성한다.

\subsection{적응 꼬리 정밀화와 후보 감사}
\label{sec:design:tail}

$N_{\rm null}=100$의 최소 PTE는 $1/101\simeq0.0099$이므로, 0.01
임계값 부근의 순위는 Monte Carlo 해상도에 민감하다. 국소 꼬리는
$PTE_{\rm global,max}\le0.03$ 또는 $Z_{\rm PS,max}>3$이면서
$B_{\rm rel}>10^{-3}$인 행을, 적분 꼬리는
$PTE_{\rm global,absint}\le0.03$이면서 $B_{\rm rel}>10^{-3}$인
행을 선택한다. 선택 후 짝지은 빔을 추가하고 모든 행을
$N_{\rm null}=\NNullTail$로 재평가한다.

국소 분기의 후보 수준 감사는 거친 그리드에서
$PTE_{\rm global,max}<0.01$ 및 $B_{\rm rel}>10^{-2}$를 만족한
국소-물리 후보 \NCandidateAudit{}건을 대상으로 수행한다. 동일한
물리 조건의 양쪽 빔을 $N_{\rm null}=1000$으로 재계산한 뒤,
백색·매끄러운 가산 잔차 대용물을
$\sigma_{\rm cal}\in\{0,10^{-4},10^{-3},10^{-2}\}$에서 적용한다.
매끄러운 잔차의 주파수 상관 길이는 0.5, 1.0, 5.0 MHz이다.

적분 분기에 대해서는 별도의 전수 감사를 수행한다. 꼬리 정밀화에서
탐색적 $10^{-3}$ 게이트를 통과한 빔-강건 적분 행
\NIntTailBeamRobustEthree{}건, 즉 짝지은 고유 쌍
\NIntAuditPairs{}개 전체를 대상으로, 국소 감사와 동일한 잔차
설계---짝지은 frozen/full PolyBeam, 조건별
$N_{\rm null}=\NNullTail$, 위와 동일한 $\sigma_{\rm cal}$ 및 잔차
모형(백색 4조건, 매끄러운 $\sigma_{\rm cal}\times\ell_\nu$
12조건)---를 적용한다. 각 (쌍, 잔차 조건)에서 양쪽 빔이 교란된
배경 하에서도 독립적으로 $PTE_{\rm global,absint}<0.01$과
$B_{\rm rel}>10^{-3}$을 유지하면 빔-강건 유지로, 한쪽 빔만
유지하면 단일 빔 유지로, 그 외에는 탈락으로 분류한다. 전체 행렬은
\NIntAuditPairs{}쌍 $\times$ \NIntAuditConds{}조건 $\times$ 2빔
$=$ \NIntAuditEvals{}회의 $N_{\rm null}=\NNullTail$ 평가로
구성되며, 잔차 실현은 국소 감사와 동일한 단일 seed를 사용한다.
$\sigma_{\rm cal}=0$ 행은 원 꼬리 정밀화 산출물과 수치적으로
일치함을 교차 검증하였다.

\subsection{합성 천정각 감도 분석}
\label{sec:design:za}

이질적인 실제 TLE 앙상블과 분리하여 기하학적 감도를 평가하기 위해
고도 550 km, 통과 방위각 $90^\circ$인 합성 궤도를 사용한다.
최고 천정각은
\begin{equation}
 z_{\rm peak}\in\{5^\circ,15^\circ,30^\circ,45^\circ,
 60^\circ,70^\circ,78^\circ\}
\end{equation}
로 변화시킨다. 5개 배경, 세 동서 방향 기선(14.6, 140.4, 207.3 m),
일곱 천정각 및 두 빔 조건을 교차하여 총 \NZAAll{}개 사례를
구성한다. full-chromatic PolyBeam은 물리 빔 조건이고,
$B\equiv1$은 빔 감쇠를 제거한 geometry-only 진단 비교군이다.
$S_{\rm ref}=300$ Jy와 smooth 형태를 고정하여 각도 효과를 분리한다.

% =============================================================================
\section{결과}
\label{sec:results}

\subsection{4-플럭스 주 그리드}
\label{sec:results:grid}

거친 \NMainRows{}행 그리드에서 엄격한 국소 행은
\NMainLocalStrict{}건, 엄격한 적분 행은 \NMainIntStrict{}건이었다.
그러나 $B_{\rm rel}>10^{-3}$을 함께 요구하면 각각
\NMainLocalPhysEthree{}건과 \NMainIntPhysEthree{}건으로 줄었다.
운영적 $10^{-2}$ 게이트에서는 국소 \NMainLocalPhysEtwo{}건,
적분 \NMainIntPhysEtwo{}건이 남았다. 짝지은 빔 기준
$10^{-3}$ 게이트에서 빔-강건 국소 쌍은
\NMainBeamRobustLocalEthree{}개, 적분 쌍은
\NMainBeamRobustIntEthree{}개였다. 참고로 시행별 PTE가 균등하게
분포한다면 \NMainRows{}행에서 $PTE<0.01$ 행의 기대 수는 약
\NMainExpectedFP{}건이다. 관측된 \NMainLocalStrict{}건과
\NMainIntStrict{}건은 이를 상회하지만, 본 PTE는
Section~\ref{sec:background}에서 정의한 조건부 민감도 지표이므로
이 비교는 규모 감각을 위한 참고값이다.

\begin{table*}[!t]
\centering
\small
\caption{플럭스별 거친 주 그리드 결과($N_{\rm null}=100$). 물리 행은
해당 분기 PTE와 $B_{\rm rel}$ 게이트를 동시에 만족한다.}
\label{tab:flux_counts}
\begin{tabular}{rrrrrrrr}
\toprule
$S_{\rm ref}$ (Jy) & 행 & L 엄격 & I 엄격 &
L+$10^{-3}$ & I+$10^{-3}$ & L+$10^{-2}$ & I+$10^{-2}$ \\
\midrule
10   & 324 & 9 & 18 & 0 & 6  & 0 & 0 \\
30   & 324 & 8 & 17 & 2 & 12 & 0 & 2 \\
300  & 324 & 6 & 16 & 4 & 16 & 2 & 10 \\
1000 & 324 & 7 & 23 & 5 & 23 & 5 & 23 \\
\bottomrule
\end{tabular}
\end{table*}

플럭스 증가에 따라 엄격 PTE 행 수가 단조롭게 증가하지는 않았지만,
절대편향 게이트를 포함한 물리 행은 고플럭스 계층에 집중되었다.
이는 대응-영가설도 주입 진폭과 함께 변하기 때문에 PTE와 효과크기가
동일한 역할을 하지 않음을 보여 준다. 형태학별
$B_{\rm rel}>10^{-3}$ 물리 행은 smooth에서 국소/적분 3/16,
lines에서 2/23, kHz-comb에서 6/18이었다.

\begin{figure}[!t]
\centering
\maybeincludegraphics[width=\linewidth]{R2_z_vs_pte_global.png}
\caption{4-플럭스 그리드의 국소 강건 z와
$PTE_{\rm global,max}$. 큰 z는 낮은 경험적 PTE 또는 큰
$B_{\rm rel}$과 일대일 대응하지 않으므로 탐색적 진단으로만 사용한다.}
\label{fig:z-pte}
\end{figure}

\begin{figure}[!t]
\centering
\maybeincludegraphics[width=\linewidth]{R3_bias_floor_sensitivity.png}
\caption{효과크기 게이트에 따른 국소·적분 물리 행 수. $B_{\rm rel}$
임계값은 QA 보고 규칙이며 물리적으로 보정된 우주론적 허용치가 아니다.}
\label{fig:bias-floor}
\end{figure}

\subsection{적분 편향을 지배하는 요인}
\label{sec:results:anova}

양의 $B_{\rm rel}$을 갖는 행에 대해 $\log_{10}B_{\rm rel}$의 Type-II
분산 분해를 수행하였다. \AnovaDropped{}개 행은 비양의 반응값으로
로그 모형에서 제외되었다. 모형의 $R^2$와 조정 $R^2$는 각각
\AnovaRtwo{}와 \AnovaAdjRtwo{}였다.

\begin{table}[!t]
\centering
\small
\caption{$\log_{10}B_{\rm rel}$ Type-II 분산 분해. $10^{-16}$
미만의 $p$-값은 수치 정밀도 한계 아래이므로 절단 표기한다.}
\label{tab:anova}
\begin{tabular}{lrr}
\toprule
요인 & partial $\eta^2$ & $p$ \\
\midrule
플럭스 & \AnovaFluxEta & $<10^{-16}$ \\
LST 층 & \AnovaLSTEta & $<10^{-16}$ \\
위성 다중성 & \AnovaMultEta & $<10^{-16}$ \\
기선 길이 & \AnovaBaseEta & $<10^{-16}$ \\
빔 모델 & \AnovaBeamEta & 0.0075 \\
형태학 & \AnovaMorphEta & 0.38 \\
빔$\times$형태학 & \AnovaInterEta & 0.94 \\
\bottomrule
\end{tabular}
\end{table}

플럭스가 가장 큰 효과를 보였고, 기술적 log--log 기울기는 플럭스
1 dex 증가당 $B_{\rm rel}$ \AnovaFluxSlope{} dex였다. LST 층,
다중성 및 기선 길이도 큰 효과를 보인 반면, 빔과 형태학의 효과크기는
작았다. 이는 빔이 중요하지 않다는 뜻이 아니다. 빔은 전체 편향의
분산을 크게 설명하지 않더라도 개별 사례를 임계값 양쪽으로 이동시켜
후보 등급을 바꿀 수 있으므로 짝지은 후보 감사가 필요하다.

\subsection{$N_{\rm null}=1000$ 병렬 꼬리 정밀화}
\label{sec:results:tail}

국소 꼬리에서는 짝지은 확장 후 \NLocalTailRows{}개 행을 재평가해
엄격 국소 \NLocalTailStrict{}건을 얻었다. 이 중
$B_{\rm rel}>10^{-3}$은 \NLocalTailPhysEthree{}건,
$10^{-2}$는 \NLocalTailPhysEtwo{}건이었고, 탐색적
$10^{-3}$ 게이트의 빔-강건 국소 행은
\NLocalTailBeamRobustEthree{}건이었다.

적분 꼬리에서는 \NIntTailRows{}개 행을 재평가해 엄격 적분
\NIntTailStrict{}건을 얻었다. $B_{\rm rel}>10^{-3}$은
\NIntTailPhysEthree{}건, $10^{-2}$는 \NIntTailPhysEtwo{}건이었으며,
선택에 사용한 $10^{-3}$ 게이트에서 빔-강건 적분 행은
\NIntTailBeamRobustEthree{}건, 짝지은 고유 쌍으로는
\NIntAuditPairs{}개였다.

\begin{table}[!t]
\centering
\small
\caption{적응 꼬리 정밀화 결과($N_{\rm null}=1000$). 빔-강건 수는
선택에 사용한 $B_{\rm rel}>10^{-3}$ 게이트 기준의 행 수이며, 적분
꼬리의 짝지은 고유 쌍 수는 \NIntAuditPairs{}개이다.}
\label{tab:tail}
\begin{tabular}{lrr}
\toprule
지표 & 국소 꼬리 & 적분 꼬리 \\
\midrule
재평가 행 & \NLocalTailRows & \NIntTailRows \\
엄격 PTE & \NLocalTailStrict & \NIntTailStrict \\
PTE+$B_{\rm rel}>10^{-3}$ & \NLocalTailPhysEthree & \NIntTailPhysEthree \\
PTE+$B_{\rm rel}>10^{-2}$ & \NLocalTailPhysEtwo & \NIntTailPhysEtwo \\
빔-강건($10^{-3}$, 행) & \NLocalTailBeamRobustEthree & \NIntTailBeamRobustEthree \\
\bottomrule
\end{tabular}
\end{table}

이 결과는 이전 2-플럭스 분석의 ``빔-강건 적분 후보 0'' 결론을
4-플럭스 결과에 그대로 적용할 수 없음을 보여 준다. 적분 꼬리의
\NIntAuditPairs{}쌍은 탐색적 $10^{-3}$ 효과크기 게이트를 통과한
QA 후보이며, 이 집합 전체에 대한 보정 잔차 감사는
Section~\ref{sec:results:intaudit}에서 보고한다.

\subsection{천정각 스윕: 기하와 빔 응답의 분리}
\label{sec:results:za}

full-chromatic PolyBeam \NZAPerBeam{}개 사례에서는
$B_{\rm rel}>10^{-3}$을 동반한 엄격 국소 및 적분 후보가 모두 0건이었다.
$B\equiv1$ 기하학 진단에서는 국소 후보 \NZANoBeamLocalPhys{}건,
적분 후보 0건이 나타났다. 효과크기만 보면 $B_{\rm rel}>10^{-2}$인
사례가 full PolyBeam에서 \NZAFullBiasEtwo{}건,
$B\equiv1$에서 \NZANoBeamBiasEtwo{}건이었지만, PTE 게이트를 함께
통과하지 못했다.

\begin{table}[!t]
\centering
\small
\caption{현재 QA 연산자로 재실행한 합성 천정각 스윕. 물리 후보는
PTE와 $B_{\rm rel}>10^{-3}$을 동시에 요구한다.}
\label{tab:za}
\begin{tabular}{lrrr}
\toprule
빔 조건 & 사례 & 국소 물리 & 적분 물리 \\
\midrule
full H2C PolyBeam & \NZAPerBeam & \NZAFullLocalPhys & \NZAFullIntPhys \\
$B\equiv1$ 진단 & \NZAPerBeam & \NZANoBeamLocalPhys & \NZANoBeamIntPhys \\
\bottomrule
\end{tabular}
\end{table}

full PolyBeam에서 가장 작은 적분 PTE는 $z_{\rm peak}=78^\circ$,
140.4 m 기선의 한 배경에서 0.0099였으나,
$PTE_{\rm global,max}=0.92$이고 $B_{\rm rel}=3.9\times10^{-5}$로
효과크기 게이트를 통과하지 못했다. 천정각 또는 지평선 근접도만으로
오염 후보를 판정할 수 없으며, 빔 결합과 배경 조건을 포함한 전체
연산자 응답이 필요하다.

\subsection{국소 후보와 보정 잔차 감사}
\label{sec:results:audit}

거친 그리드에서 $B_{\rm rel}>10^{-2}$을 동반한 국소-물리 후보
\NCandidateAudit{}건을 $N_{\rm null}=1000$과 짝지은 빔으로
재평가하였다. \NCandidateLocalRetained{}건은 엄격한 국소
유의성을 유지했지만 적분 분기는 비유의했고, 1건은
$PTE_{\rm global,max}=\NCandidateDroppedPTE$로 0.01 임계값을
벗어났다. 이 제한된 국소 후보 집합에서 빔 민감 후보와 빔-강건 적분
오염 후보는 각각 \NCandidateBeamSensitive{}건과
\NCandidateBeamRobustContam{}건이었다.

동일한 7개 후보에 백색 및 매끄러운 보정 잔차를 적용한 모든 조건에서
분류는 6개 엄격 국소/적분 비유의 사례와 1개 국소 임계값 이탈 사례로
유지되었다. 어떤 잔차 조건도 후보를 빔 민감 또는 빔-강건 적분 오염
등급으로 승격하지 않았다. 이 감사는 국소 후보의 과대해석을
제한하며, 적분 꼬리 후보의 전수 감사는
Section~\ref{sec:results:intaudit}에서 별도로 수행한다.

\subsection{적분 꼬리 빔-강건 후보의 전수 잔차 감사}
\label{sec:results:intaudit}

꼬리 정밀화의 빔-강건 적분 집합은 \NIntTailBeamRobustEthree{}행,
짝지은 고유 쌍 \NIntAuditPairs{}개로 확정되었다. 이 중
\NIntAuditPairsEtwo{}쌍은 양쪽 빔에서 운영적
$B_{\rm rel}>10^{-2}$ 게이트도 통과했고, 나머지
\NIntAuditPairsMid{}쌍(30 Jy, smooth 및 lines)은
$10^{-3}$--$10^{-2}$ 사이의 $B_{\rm rel}$을 보였다. 후보는 두
기선--LST 셀에 집중되었다. 전형 층의 한 셀에서
\NIntAuditPairsTypical{}쌍(세 형태학 모두, 1000 Jy, 다중 위성),
스트레스 층의 한 셀에서 \NIntAuditPairsStress{}쌍(30--1000 Jy,
다중 위성)이 나타났으며, 단일 위성 또는 10 Jy 조건의 쌍은 없었다.

\NIntAuditPairs{}쌍 $\times$ \NIntAuditConds{}개 잔차 조건의 전체
\NIntAuditCells{}개 (쌍, 조건) 셀에서, 어떤 조건도 단일 쌍을 단일
빔 유지 또는 탈락 등급으로 강등시키지 않았다. 전 행렬에 걸쳐 양쪽
빔의 $PTE_{\rm global,absint}$ 최댓값은 \IntAuditMaxPTE{}로 0.01
게이트 대비 약 30\%의 여유를 유지했고, $B_{\rm rel}$ 최솟값은
$\IntAuditMinBrel$로 $10^{-3}$ 게이트의 약 1.9배였다. 기준선에서
가장 작은 여유를 보인 두 30 Jy 쌍도 전체 잔차 행렬에서
$B_{\rm rel}$이 $[1.9,2.4]\times10^{-3}$ 범위 안에서만 움직였으며,
$10^{-3}$ 바닥으로 떨어지지도 $10^{-2}$ 계층으로 올라서지도
않았다. $10^{-2}$ 계층 \NIntAuditPairsEtwo{}쌍은 최강 교란
조건($\sigma_{\rm cal}=10^{-2}$)에서도 양쪽 빔 모두
$B_{\rm rel}>10^{-2}$를 유지했다(최소 \IntAuditMinBrelEtwo{}).

\begin{table}[!t]
\centering
\small
\caption{적분 꼬리 빔-강건 쌍의 전수 보정 잔차 감사 요약. 마진은
\NIntAuditCells{}개 (쌍, 조건) 셀 전체와 양쪽 빔에 걸친 극값이다.}
\label{tab:intaudit}
\begin{tabularx}{\linewidth}{Xr}
\toprule
지표 & 값 \\
\midrule
감사 대상(행 / 짝지은 쌍) & \NIntTailBeamRobustEthree{} / \NIntAuditPairs \\
잔차 조건(백색 + 매끄러운) & \NIntAuditConds \\
$N_{\rm null}=\NNullTail$ 평가 횟수 & \NIntAuditEvals \\
빔-강건 유지 (쌍, 조건) 셀 & \NIntAuditCells{} / \NIntAuditCells \\
최대 $PTE_{\rm global,absint}$ (게이트 $<0.01$) & \IntAuditMaxPTE \\
최소 $B_{\rm rel}$ (게이트 $>10^{-3}$) & $\IntAuditMinBrel$ \\
$B_{\rm rel}>10^{-2}$ 계층 유지 쌍 & \NIntAuditPairsEtwo{} / \NIntAuditPairsEtwo \\
\bottomrule
\end{tabularx}
\end{table}

\begin{figure}[!t]
\centering
\maybeincludegraphics[width=\linewidth]{R5_residual_audit_margins.png}
\caption{적분 꼬리 빔-강건 \NIntAuditPairs{}쌍의 잔차 감사 마진.
쌍별로 \NIntAuditConds{}개 잔차 조건과 양쪽 빔에 걸친
$B_{\rm rel}$의 범위(위)와 $PTE_{\rm global,absint}$의
최댓값(아래)을 표시하며, 점선은 $10^{-3}$·$10^{-2}$ 효과크기
게이트와 0.01 PTE 게이트이다. 어떤 쌍도 어떤 조건에서도 게이트를
이탈하지 않는다.}
\label{fig:intaudit}
\end{figure}

따라서 Table~\ref{tab:hierarchy}의 최종 단계 기준으로, 이
\NIntAuditPairs{}쌍은 오염 후보 승격 조건을 만족한다. 강조하건대
이는 주입된 스트레스 신호에 대한 결과이다. 승격이 뜻하는 바는
빔-강건 적분 분류가 특정 (무교란) 보정 실현의 산물이 아니라는
것, 즉 판정 계층이 잔차 교란 하에서도 안정적으로 동일한 후보
집합을 지목한다는 것이며, 실제 21-cm 자료에서 오염이 검출되었다는
뜻이 아니다.

\FloatBarrier

% =============================================================================
\section{논의}
\label{sec:discussion}

첫째, 국소성과 적분성은 실제로 다른 후보 집합을 만든다. 거친
그리드와 정밀화 모두에서 적분 분기 후보가 국소 분기 후보보다 많았고,
국소 후보 7건의 정밀 감사에서는 적분 유의성이 유지되지 않은 반면,
적분 꼬리의 빔-강건 쌍은 전수 잔차 감사를 통과했다. 따라서 국소
경보를 윈도우 전체의 오염으로 읽거나, 반대로 국소 분기를 적분
분기의 선행 관문으로 사용하는 것은 타당하지 않다.

둘째, 4단계 플럭스 축은 PTE와 효과크기의 역할을 분리한다. PTE
후보 수는 플럭스에 대해 단조롭지 않았지만, $B_{\rm rel}$을 포함한
물리 행과 Type-II 분산 분해는 플럭스가 적분 편향의 지배적 요인임을
보였다. 10 및 30 Jy는 문헌 규모의 anchor이고, 300 및 1000 Jy는
실제 발생률을 뜻하지 않는 상한 스트레스이므로, 모든 결과는 플럭스
계층별로 분리해 해석해야 한다.

셋째, 빔은 전역 분산을 크게 설명하지 않더라도 후보 판정을 바꿀 수
있다. full/frozen PolyBeam은 동일한 CST 기반 공간 패턴을 공유하지만
주파수 색수차 처리에서 다르다. 짝지은 빔 게이트는 빔 강건성만으로
물리적 실재를 증명하려는 것이 아니라, 한 연산자 가정에만 의존하는
후보를 식별하는 반증 단계이다. 합성 천정각 스윕에서 $B\equiv1$
진단이 full PolyBeam보다 훨씬 많은 큰 $B_{\rm rel}$ 사례를 만든
결과도 빔 감쇠를 제거한 기하학 상한을 물리적 후보로 해석해서는
안 됨을 보여 준다.

넷째, 적분 꼬리의 빔-강건 후보는 보정 잔차 감사를 전수
통과하였다. 빔-강건 적분 분류가 특정 보정 실현의 산물일 가능성은
\NIntAuditCells{}개 셀 전수 감사에서 배제되었다. 어떤 잔차
조건도 단일 쌍을 강등시키지 않았고, PTE와 $B_{\rm rel}$의 최악
마진은 게이트 대비 각각 약 30\%와 1.9배였다. 이로써
Table~\ref{tab:hierarchy}의 최종 단계가 처음으로 발동하며, 이
결과는 두 방향으로 읽어야 한다. 방법론적으로는 적분 분기가 빔
연산자 교체와 보정 잔차 실현 모두에 강건함을 보이는 연산자
검증이다. 물리적으로는 승격 후보가 coherent-default 다중 위성
조건과 특정 배경 셀에 국한되며, 문헌 anchor인 30 Jy에서
\NIntAuditPairsMid{}쌍, 상한 스트레스 300--1000 Jy에서
\NIntAuditPairsEtwo{}쌍이 나타났다. 30 Jy 쌍의 존재는 문헌 규모의
UEMR도 불리한 배경과 다중 위성 결맞음 하에서는 적분 QA를 발동시킬
수 있음을 시사하지만, coherent-default 합산이 상한 스트레스
구성이므로 발생률 주장으로 확장할 수 없다. 아울러 운영적
$10^{-2}$ 계층으로 게이트를 조이더라도 \NIntAuditPairsEtwo{}쌍이
잔차-안정적으로 남으므로, 본 감사의 강건성 결론은 게이트 선택에
민감하지 않다.

이 프레임워크는 관측 연구와 생산 수준 우주론 분석 사이에 위치한다.
Starlink UEMR의 존재와 스펙트럼 특성을 확립한 관측 연구
\citep{divruno2023,bassa2024}가 입력 현상을 제공하고, HERA의
밴드파워·jackknife·종단간 검증이 최종 과학 산출물을 검증한다면,
본 연구는 처리 후 잔차의 통계적 초과를 후속 검토 등급으로 분류하는
pre-production QA 규칙을 제공한다
\citep{abdurashidova2022,wilensky2023jackknife}.

\subsection{제한점}

결론은 2022년 10월 8일 HERA H6C 파생 배경을 2026년 1월 1일 TLE
평가 창으로 재매핑한 구성에 조건부이다. 정확한 HERA 파일 식별자,
TLE 공급자·획득 시각·필터링 규칙 및 각 TLE epoch 차이를 재현
패키지에 명시해야 한다. 전체 카탈로그 사용은 임의의 1200개
부분집합 편향을 제거하지만, 위성별 실제 방출 상태와 궤도 오차를
제거하지 않는다.

대응-영가설은 UEMR 부재 영가설이 아니다. 단일 위성에서는 주로
배경--위성 교차항 위상을, 다중 위성에서는 위성 간 상대 위상을
변경한다. coherent-default 다중 위성 합산도 실제 방출원 위상 동기를
예측하는 모형이 아니라 상한 스트레스 구성이다. 다중 위성 상대 위상
앙상블과 Doppler drift의 체계적 감사가 후속 단계에 필요하다.

$B_{\rm rel}$ 임계값은 실제 우주론적 오차 예산으로 보정되지 않았다.
생산 수준 적용에는 $\mathrm{mK}^2$ 밴드파워 정규화, 전체 열잡음
공분산, 중복도 가중치, 다중 LST 결합, 신호 손실 및 계통오차 예산이
필요하다
\citep{cheng2018,aguirre2021validation,abdurashidova2022}.
보정 감사는 가산적 백색·매끄러운 잔차에 한정되어 곱셈형 이득,
대역통과 리플, 편광 누출 및 중복 보정 잔차를 포함하지 않으며
\citep{barry2016,ewallwice2017,dillon2020redcal,kim2023beamvar},
잔차 실현도 단일 seed에 기반한다. 다만 관측된 최악 마진이 게이트
대비 충분히 크므로, 잔차 실현 교체가 본 감사의 분류를 바꿀
가능성은 제한적이다.

% =============================================================================
\section{결론}
\label{sec:conclusion}

본 연구는 전체 \NTLERecords{}개 TLE, 4단계 플럭스, HERA H2C CST
PolyBeam 및 병렬 국소/적분 PTE를 결합한 재현 가능한 지연 스펙트럼
QA 프레임워크를 제시하였다. \NMainRows{}행 주 그리드에서 효과크기를
포함한 후보 수는 플럭스에 따라 증가했고, $\log_{10}B_{\rm rel}$은
플럭스, LST 층, 위성 다중성 및 기선 길이가 지배하였다. 빔과 형태학의
전역 효과크기는 작았지만, 빔 교체는 개별 후보의 임계값 통과 여부를
바꾸므로 짝지은 연산자 감사가 필요했다.

$N_{\rm null}=1000$ 정밀화에서는 탐색적 $B_{\rm rel}>10^{-3}$에서
빔-강건 국소 QA 행 \NLocalTailBeamRobustEthree{}건과 적분 QA 행
\NIntTailBeamRobustEthree{}건(짝지은 \NIntAuditPairs{}쌍)이
확인되었고, 적분 \NIntAuditPairs{}쌍 전체는 \NIntAuditConds{}개
보정 잔차 조건의 전수 감사에서 빔-강건 분류를 유지하여 본 QA
계층의 오염 후보 승격 기준을 만족하였다. 따라서 현 단계의 정확한
결론은 ``후보가 없다''가 아니라, ``주입 스트레스 하에서
국소·적분·효과크기·빔·잔차 게이트를 모두 통과하는 잘 정의된 후보
집합이 존재하며, 이 집합은 다중 위성·고플럭스 구성과 특정 배경
셀에 집중된다''이다. 반면 합성 천정각 스윕에서는 full PolyBeam
조건의 물리적 후보가 없었고, 국소 후보 7건의 독립 잔차 감사에서도
빔-강건 적분 오염 후보가 생성되지 않았다.

이 구조는 단일 PTE, 단일 효과크기 또는 단일 빔 가정의 초과를 실제
21-cm 오염으로 과대해석하지 않도록 하는 동시에, 모든 게이트를
통과한 후보를 재현 가능하게 지목한다. 향후 실제 동시 관측과 생산
수준 파이프라인에서는 승격 후보의 다중 위성 상대 위상·Doppler
강건성과 우주론 단위의 오차 예산을 추가함으로써 본 QA 계층을
완성할 수 있다.

\section*{감사의 글}
저자들은 HERA 파생 가시성 자료와 관측·처리 기반을 공개한 HERA
Collaboration과 참여 연구진에 감사한다. 본문의 분석과 해석은
저자들의 책임이며 HERA Collaboration의 공식 견해를 나타내지 않는다.

\section*{데이터 및 코드 가용성}
분석 코드, 활성 설정, 요약 산출물 및 그림 생성 코드는
\url{https://github.com/kimeujin03-droid/starlink_uemr}에서 제공한다.
4-플럭스 핵심 산출물은 \texttt{coverage\_robustness\_trials\_fourflux.csv},
국소·적분 꼬리 정밀화 CSV, ZA 스윕 CSV 및 후보 보정 잔차 감사 CSV를
포함한다. 적분 꼬리 전수 감사 산출물로
\texttt{beam\_robust\_absint\_pairs\_fourflux\_1e3.csv}(후보 22행/11쌍),
\texttt{beam\_robust\_residual\_audit\_fourflux.csv}(176행 결과 행렬),
동 요약 CSV 및 \texttt{run\_beam\_robust\_residual\_audit.py}를
제공한다. 원시 HERA UVH5, 대용량 null 앙상블과 외부 TLE 입력은
재배포하지 않으며, 전체 재실행에는 원자료와 정확한 TLE 출처가
필요하다.

\bibliographystyle{elsarticle-num}
\bibliography{starlink_uemr_refs_complete}

\end{document}